Wind energy is a central pillar of the global transition to renewable
power, and its continued growth depends on understanding how
turbines interact with the turbulent atmospheric boundary layer (ABL)
and with each other's wakes~\cite{veers2019}. Wake losses in large wind
farms reduce total power output by 10--20\%~\cite{barthelmie2009}, and
mitigating them requires simulating tens of turbines embedded in
kilometers of turbulent atmosphere over hours of physical time.
Large-eddy simulation (LES) resolves these energy-carrying turbulent
motions directly and has become the de facto standard for wind-farm
aerodynamics~\cite{calaf2010, stevens2017}. However, it comes at a steep
computational price when coupling hundreds of millions of grid points to
the long integration times these flows demand.

LESGO~\cite{lesgo-web,bouzeid2005} is a widely used open-source pseudo-spectral
LES solver for ABL and wind-farm flows.  The code has been used for many
different kinds of simulations including wind-farm, rough-surface, canopy,
urban, and wall-bounded flows~\cite{lesgo-publications}. \fix{you need to cite dozens of publications here} Recent work continues
to use LESGO for large wind-farm simulations and public
datasets~\cite{zhu2025jhtdbwind}. \fix{what does use LESGO for public datasets mean?} Its numerical formulation combines Fourier
colocation in the horizontal directions, a scale-dependent Lagrangian dynamic
subgrid-scale (SGS) model~\cite{bouzeid2005}, and actuator-line turbine
coupling~\cite{sorensen2002}. This formulation carries little numerical
dissipation and delivers high accuracy per grid point. Its CPU implementation,
however, saturates node memory bandwidth even at tens of cores and exhausts its
slab decomposition at modest process counts (Section~\ref{sec:bg-impl}), severely
constraining production-scale runs.  Because the scientific fidelity of
turbulent wake simulations relies on capturing long physical integration times,
this performance ceiling directly limits accessible physical time scales.
Porting the solver to GPUs provides a natural path to overcome these
scalability barriers.  \rev{Figure~\ref{fig:wf60-velocity} illustrates the
interacting turbulent wakes in the 60-turbine configuration used for our
production benchmark.} \fix{this statement looks abrupt here, maybe move it to where you talk about dozens of simulations}

Porting large Fortran solvers to GPUs with compiler directives such as
OpenACC is by now well established across weather and climate
models~\cite{lapillonne2014,govett2017nim,norman2015,giorgetta2022,esclapez2025dales},
CFD and turbulence
solvers~\cite{xia2015,otero2019nek,devanna2023uranos}, and production
physics codes in other domains~\cite{caplan2019mas}. By
annotating existing loops instead of rewriting them in a GPU-specific programming
language, this approach keeps a single source tree, which lowers the 
barrier to GPU adoption for large and physics-validated code bases, and lets their user communities maintain the source directly. 
It is especially effective at offloading
compute-intensive kernels, where a self-contained block of arithmetic
moves to the GPU wholesale and improves kernel throughput. For
LESGO, however, the factor that governs performance is
not kernel throughput but data movement.

Pseudo-spectral multiphysics LES has recurring host--device
\emph{coupling points} -- places where host-side code must read or write
device-resident state. They vary widely in how much of the field they
touch and how often, so no single blanket data-placement policy fits
them all. The two data strategies common to directive-based ports both
encounter severe limitations here. (1)~Array mirroring, whether 
wholesale~\cite{esclapez2025dales} or selective~\cite{giorgetta2022}, 
synchronizes data at module boundaries. However, servicing coupling 
points over PCIe still leaves the GPU largely idle.
Moreover, blanket duplication wastes device memory on unconsumed arrays,
restricting per-GPU grid capacity.
(2)~Delegating placement to CUDA managed memory lets the driver migrate
pages on demand, a route that minimizes porting effort and underpins
Fortran standard-parallelism offloading~\cite{caplan2023dc}, yet it
hides traffic from the source code and incurs page-migration overheads 
across a range of GPU workloads~\cite{chien2019um}.
Thus, achieving high performance for LESGO requires developing a
fine-grained data residency strategy tailored to its distinct
host--device coupling points.

In this paper, we present our approach to porting the Fortran solver in LESGO
to GPGPUs. The port preserves the solver's validated physics and brings
production-scale wind-farm LES within practical reach. We begin with an analysis of the 
data movement that measures two established placement strategies \fix{what two strategies?}
on our own code and traces the bottleneck to host--device coupling points.
Guided by this analysis, we analyze all coupling points across the solver,
classify them by host consumer read footprints, and apply four treatments
(gate, slice, relocate, and overlap). \fix{I am not sure what treatements means} Beyond data-residency optimizations, we
apply three other optimizations to overcome remaining structural bottlenecks. We pipeline
independent systems across a serialized rank chain to accelerate the pressure
equation's distributed tridiagonal solve; batch many small work items to
mitigate kernel launch overhead; and hide residual host-side airfoil physics
behind device execution through CPU--GPU heterogeneous execution.

We implement our GPU port in OpenACC using the NVHPC compiler,
and evaluate it on NVIDIA A100 GPUs and H200 GPUs. 
This work makes the following contributions:
\begin{itemize}
\item We present the first ever GPU port of LESGO, covering its full
physics pipeline. To our knowledge, this is the first directive-based 
GPU LES with actuator-line wind-farm coupling at this scale. Our code 
will be open-sourced upon acceptance.
\item We analyze and classify all host--device coupling points by
their host consumer footprints, enforcing minimal data movement through
four distinct treatments (gate, slice, relocate, and overlap).
\item We demonstrate three optimizations on the
solver's coupling-heavy modules: chunk pipelining across serial rank
chains to hide inter-rank communication delays, batching many small
work items to mitigate kernel launch overhead, and CPU--GPU
heterogeneous execution to overlap residual host physics with the GPU
backlog.
\item We evaluate the GPU version at production scale on a 604M-cell, 60-turbine 
wind farm, achieving a 40$\times$ speedup over the fastest CPU baseline.
We also demonstrate scaling up to 128 GPUs on a 1.36 billion-cell grid. 
\end{itemize}
